% Fandol has no OpenType CJK script tag. xeCJK requests Script=CJK, fontspec
% warns, and then falls back to the default script. CustomDefault is
% fontspec's quiet DFLT/latn fallback and selects that same result.
\AddToHook{package/xeCJK/after}{%
  \defaultCJKfontfeatures{Script=CustomDefault}%
}

\documentclass[12pt,fontset=fandol]{ctexart}

% 示例信息与测量值均为虚构内容，提交前请替换。
\newcommand{\reporttitle}{一阶阻容电路阶跃响应实验}
\newcommand{\reportsubtitle}{RC Step Response}
\newcommand{\coursename}{电路实验}
\newcommand{\studentname}{张同学}
\newcommand{\studentid}{20260000000}
\newcommand{\instructorname}{李老师}
\newcommand{\experimentdate}{2026年4月18日}

\usepackage[a4paper,top=2.5cm,bottom=2.5cm,left=2.6cm,right=2.6cm,headheight=16pt]{geometry}
\usepackage{setspace}
\usepackage{fancyhdr}
\usepackage{graphicx}
\usepackage{booktabs}
\usepackage{array}
\usepackage{amsmath,amssymb}
\usepackage{siunitx}
\usepackage{xcolor}
\usepackage{circuitikz}
\usepackage{pgfplots}
\usetikzlibrary{arrows.meta}
\pgfplotsset{compat=1.18}
\usepackage{csquotes}
\usepackage[backend=biber,style=gb7714-2015]{biblatex} % chktex 8
\addbibresource{references.bib}
\usepackage{hyperref}
\usepackage{bookmark}

\hypersetup{
  % Page numbers restart after the title page. Unique anchors avoid a
  % duplicate @page destination in xdvipdfmx.
  hypertexnames = false,
  colorlinks = true,
  linkcolor = blue!45!black,
  citecolor = green!35!black,
  urlcolor = blue!55!black,
  pdftitle = {\reporttitle},
  pdfauthor = {\studentname}
}

\setstretch{1.3}
\setlength{\parindent}{2em}
\setlength{\parskip}{0.2em}
\sisetup{
  mode = match,
  per-mode = symbol
}

\pagestyle{fancy}
\fancyhf{}
\fancyhead[L]{\coursename}
\fancyhead[R]{\studentname}
\fancyfoot[C]{\thepage}

\newcommand{\metadatarow}[2]{%
  \makebox[4.5em][l]{#1} & #2\\[0.85em]%
}


\begin{document}

\begin{titlepage}
  \centering
  \vspace*{1.6cm}
  {\zihao{2}\bfseries \reporttitle\par}
  \vspace{0.45cm}
  {\Large \reportsubtitle\par}
  \vspace{2.4cm}
  \renewcommand{\arraystretch}{1.15}
  \begin{tabular}{@{}rl@{}}
    \metadatarow{课程名称}{\coursename}
    \metadatarow{姓名}{\studentname}
    \metadatarow{学号}{\studentid}
    \metadatarow{指导教师}{\instructorname}
    \metadatarow{实验日期}{\experimentdate}
  \end{tabular}
  \vfill
  {\small 本模板为非官方实验报告示例。表中测量点和曲线偏移都是排版用的虚构数据。\par}
\end{titlepage}

\section*{实验目的}
观察一阶阻容电路的阶跃响应，练习时间常数的计算、有效数字和不确定度表述\cite{nilsson2015electric,bipm2019si}。

\section{原理}
图~\ref{fig:rc} 中的电容电压在零状态阶跃输入下满足
\begin{equation}
  v_{C}(t)=V_{s}\left(1-\mathrm{e}^{-t/\tau}\right),
  \qquad \tau=RC.
  \label{eq:step}
\end{equation}
示例参数取 $R=\qty{10.0}{\kilo\ohm}$、$C=\qty{10.0}{\micro\farad}$，因此 $\tau=\qty{0.100}{\second}$。阶跃幅值 $V_{s}=\qty{5.00}{\volt}$。

\begin{figure}[htbp]
  \centering
  \begin{circuitikz}[european, >=Stealth]
    \draw (0,0) to[V=$V_{s}$] (0,2.6)
          to[R=$R$] (3.4,2.6)
          to[C=$C$] (3.4,0)
          to[short] (0,0);
  \end{circuitikz}
  \caption{一阶阻容实验电路}\label{fig:rc}
\end{figure}

\section{仪器与器件}
\begin{table}[htbp]
  \centering
  \caption{示例仪器清单}\label{tab:apparatus}
  \begin{tabular}{lll}
    \toprule
    名称 & 示例型号 & 本次用途 \\
    \midrule
    直流电源 & DEMO-PS & 阶跃输入 \\
    数字示波器 & DEMO-SCOPE & 记录 $v_{C}(t)$ \\
    电阻 & \qty{10.0}{\kilo\ohm} & 时间常数 \\
    电容 & \qty{10.0}{\micro\farad} & 时间常数 \\
    \bottomrule
  \end{tabular}
\end{table}

\section{数据与曲线}
表~\ref{tab:data} 和图~\ref{fig:curve} 只演示排版。示意测量相对式~\eqref{eq:step} 有意偏离，不能当作真实实验结论。

\begin{table}[htbp]
  \centering
  \caption{虚构采样点}\label{tab:data}
  \begin{tabular}{
      S[table-format=1.2]
      S[table-format=1.2]
      S[table-format=1.2]
    }
    \toprule
    {时间/\unit{\second}} & {模型/\unit{\volt}} & {示意测量/\unit{\volt}} \\
    \midrule
    0.05 & 1.97 & 1.90 \\
    0.10 & 3.16 & 3.05 \\
    0.20 & 4.32 & 4.20 \\
    0.40 & 4.91 & 4.85 \\
    \bottomrule
  \end{tabular}
\end{table}

\begin{figure}[htbp]
  \centering
  \begin{tikzpicture}
    \begin{axis}[
        width=0.82\textwidth,
        height=6.4cm,
        xlabel={$t/\mathrm{s}$},
        ylabel={$v_{C}/\mathrm{V}$},
        xmin=0, xmax=0.5,
        ymin=0, ymax=5.4,
        legend pos=south east,
        grid=both,
        tick label style={font=\small},
        label style={font=\small}
      ]
      \addplot[blue, thick, domain=0:0.5, samples=80] {5*(1-exp(-x/0.1))};
      \addlegendentry{式~\eqref{eq:step}}
      \addplot[only marks, mark=*, mark size=1.7pt] coordinates {
        (0.05,1.90) (0.10,3.05) (0.20,4.20) (0.40,4.85)
      };
      \addlegendentry{示意测量}
    \end{axis}
  \end{tikzpicture}
  \caption{阶跃响应的示意曲线}\label{fig:curve}
\end{figure}

\section{不确定度}
示波器电压读数的示意不确定度取 \qty{0.05}{\volt}，覆盖仪器分辨率和读数散布。正式报告应给出评定方法、包含因子和单位\cite{bipm2019si}。本示例没有做重复测量，因此不报告实验标准偏差。

\section{结论}
在所给示例参数下，模型时间常数为 \qty{0.100}{\second}。提交前应替换为实测电阻、电容、原始波形和不确定度，并删除本段中的虚构采样点。

\printbibliography[heading=bibintoc]

\end{document}
